\begin{corollaire}
\label{IX.6.6}
Sea~$f\colon X\to S$ un morfismo propio de preesquemas localmente
noetherianos tal que~$f_*(\cal{O}_X)=\cal{O}_S$, y sea~$X'$ un preesquema
propio sobre~$X$. Consideremos la factorización de \emph{Stein}
~$X'\to S'\to S$ de~$X'\to S$ y el morfismo canónico
~$X'\to X\times_S S'$. Sean~$s$ un punto de~$S$ y~$s'$ un punto de~$S'$
situado sobre~$s$, correspondiente a una componente conexa~$F_1'$ de la
fibra~$X'_s$ de~$X'$ en~$s$. Para que el morfismo
~$X'\to X\times_S S'$ sea un isomorfismo sobre un entorno abierto~$U'$ de
~$s'$ étale sobre~$S$, es necesario y suficiente que~$X'$ sea étale sobre
~$X$ en los puntos de~$F_1'$ y que~$F_1'$ sea un recubrimiento
geométricamente trivial de la fibra~$F=X_s$.
\end{corollaire}

La necesidad sigue siendo trivial, por lo que resta probar la suficiencia.
La conclusión significa también que a) el morfismo deducido de
~$X'\to X\times_S S'$ por el cambio de base
~$\Spec(\widehat{\cal{O}}_{s'})\to S'$ es un isomorfismo; b)~$S'$ es étale
sobre~$S$ en~$s'$, es decir,~$\widehat{\cal{O}}_{s'}$ es étale sobre
~$\widehat{\cal{O}}_s$. Bajo esta forma, se ve que la conclusión es
invariante por el cambio de base~$\Spec(\widehat{\cal{O}}_s)\to S$. Como las
hipótesis son también estables por este cambio de base, podemos suponer que
~$S$ es el espectro de un anillo local noetheriano completo. Podemos además
suponer evidentemente que~$X'$ es conexo, lo cual implica aquí que
~$S'=\Spec(\cal{O}_{s'})$ y~$F'=F_1'$.

Como
\marginpar{257}
el conjunto de los puntos de~$X'$ en los que~$X'$ es étale sobre~$X$ es
abierto y contiene la fibra~$X'_{s'}=F'$, y como~$X'$ es propio sobre~$S$, se
sigue que~$X'$ es étale sobre~$X$. Puesto que induce sobre~$F=X_s$ un
recubrimiento étale isomorfo a un~$F\otimes_{\kres(s)}L$, donde~$L$ es étale
sobre~$\kres(s)$, resulta de~\Ref{IX.1.10} que es isomorfo a un
recubrimiento de la forma~$X\times_S T$, con~$T$ étale sobre~$S$ (N.B.: aquí
también basta utilizar que el funtor de~\Ref{IX.1.10} es plenamente fiel,
lo cual resulta de que todo isomorfismo formal de haces coherentes sobre~$X$
procede de un isomorfismo de esos haces). Por tanto, si~$T$ está definido por
el álgebra~$B$ finita sobre~$A$,~$X'$ se identifica con el espectro del
álgebra~$\cal{O}_X\otimes_A B$ sobre~$X$, de donde resulta inmediatamente,
puesto que~$f_*(\cal{O}_X)=\cal{O}_S$, que~$h_*(\cal{O}_{X'})$ está definido
por~$B$. Así, el homomorfismo canónico~$X'\to X\times_S S'$ no es otro que el
isomorfismo considerado~$X'\isomto X\times_S T$. Esto termina la
demostración.

\begin{corollaire}
\label{IX.6.7}
Bajo las condiciones de~\Ref{IX.6.5}, para que existan un preesquema~$S'$
étale sobre~$S$ y un~$X$-isomorfismo~$X'\isomto X\times_S S'$, es
necesario y suficiente que~$X'$ sea étale sobre~$X$ y que, para todo
~$s\in S$,
\ifthenelse{\boolean{orig}}{$X'_{s'}$}{$X'_s$}
sea un recubrimiento geométricamente trivial de~$X_s$.
\end{corollaire}

En efecto, si esto ocurre,~$X'$ es unión de abiertos~$X_i'$ isomorfos a
imágenes inversas de~$S_i'$ étale sobre~$S$. Se ve entonces fácilmente que
estos~$S_i'$ se pegan en un~$S'$ étale sobre~$S$ y que se obtiene un
isomorfismo~$X'\isomto X\times_S S'$. Puede decirse, por ejemplo, que los
~$X_i'$ están provistos de datos de descenso relativos a~$X\to S$, que se
pegan necesariamente en un dato de descenso sobre~$X'$ entero relativo a
~$X\to S$. Y, como este último es efectivo sobre los~$X_i'$, se sigue
fácilmente (gracias a un sorites omitido en el número~\Ref{IX.4}) que es
efectivo. También puede enunciarse~\Ref{IX.6.7} en la forma siguiente:

\begin{corollaire}
\label{IX.6.8}
Sea~$f\colon X\to S$ un morfismo propio y sobreyectivo de presentación
finita, con fibras geométricamente conexas. Entonces~$f$ es un morfismo de
descenso efectivo para la categoría fibrada de los preesquemas étale y
finitos sobre otros. El funtor~$S'\mto X\times_S S'$ induce una equivalencia
de la categoría de los preesquemas étale y finitos sobre~$S$ con la
categoría de los preesquemas
\marginpar{258}
étale y finitos sobre~$X$ que inducen sobre cada fibra~$X_s$ un
recubrimiento geométricamente trivial.
\end{corollaire}

\begin{remarque}
\label{IX.6.9}
Sea~$f\colon X\to S$ un morfismo propio y sobreyectivo, con
\ifthenelse{\boolean{orig}}{$Y$}{$S$}
localmente noetheriano. Entonces~$f$ se factoriza en un morfismo~$X\to S'$
que satisface la hipótesis de~\Ref{IX.6.8} y un morfismo finito y sobreyectivo
~$S'\to S$ al que se aplica~\Ref{IX.4.7}; por tanto, se expresa como
composición de dos
morfismos que son \emph{morfismos de descenso efectivo universales} para la
categoría fibrada de los preesquemas étale y finitos sobre otros. Puede
concluirse que el propio~$f$ es un morfismo de descenso efectivo universal
para la categoría fibrada considerada. Se recupera así~\Ref{IX.4.12} por un
método diferente.
\end{remarque}

\begin{remarque}
\label{IX.6.10}
La conclusión de~\Ref{IX.6.7} deja de ser válida si se sustituye la hipótesis
de que~$f$ es propio por la siguiente:~$X$ es de tipo finito sobre~$S$ y
admite una sección sobre~$S$ (por tanto,~$f$ es universalmente submersivo y
un morfismo de descenso para la categoría fibrada de los preesquemas étale
sobre otros), incluso cuando~$S$ es el espectro de un anillo de valoración
discreta y~$X'$ es un recubrimiento étale de~$X$. Para verlo, se parte de un
~$Z$ propio sobre~$S$, cuya fibra genérica
es una curva racional no singular y cuya fibra especial~$Z_0$ consta de dos
rectas secantes. Por ejemplo, si~$t$ es un uniformizante del anillo de
valoración~$A$, se toma el subesquema cerrado~$Z$ de~$\PP_A^2$ definido por
la ecuación homogénea~$x^2+y^2+tz^2=0$ (coordenadas homogéneas~$x$,~$y$,
~$z$). Se toma por~$X$ el complementario del punto singular~$a$ de~$Z_0$ en
la unión~$Z\cup\PP_k^2$. Las fibras de~$X$ son~$\PP_k^1$ y
~$\PP_k^2-a$; son, por tanto, geométricamente simplemente conexas (es decir,
todo recubrimiento étale de una de estas fibras es geométricamente trivial).
Sin embargo,
\marginpar{259}
procediendo como en el número~\Ref{IX.4}, se construyen fácilmente
recubrimientos étale de~$X$ que no proceden de recubrimientos étale de~$S$,
pegando recubrimientos triviales de~$Z-a$ y de~$\PP_k^2-a$. En cambio, es
posible que la conclusión de~\Ref{IX.6.7} subsista si se sustituye la
hipótesis de que el morfismo sea propio por la de que~$X$ sea
\emph{universalmente abierto} y
de presentación finita sobre~$S$\kern1pt\footnote{Esto ya está demostrado,
siendo~\ifthenelse{\boolean{orig}}{$g$}{$f$} solamente universalmente abierto
y sobreyectivo; \cf SGA~4 XV~1.15.}.
Esto es cierto al menos si se supone que las fibras de~$X$ sobre~$S$ son
geométricamente irreducibles, y no solo geométricamente conexas. Señalemos
únicamente que, en esta cuestión, podemos reducirnos al caso en que~$S$ es
el espectro de un anillo de valoración discreta completo con cuerpo residual
algebraicamente cerrado.
\end{remarque}

La interpretación de~\Ref{IX.6.7} en términos del grupo fundamental es la
siguiente:

\begin{corollaire}
\label{IX.6.11}
Sea~$f\colon X\to S$ un morfismo propio y sobreyectivo de presentación
finita, con fibras geométricamente conexas. Supongamos que~$X$, y por tanto
~$S$, es conexo. Sean~$a$ un punto geométrico de~$X$,~$b$ su imagen en~$S$ y,
para todo~$s\in S$, elijamos una clausura algebraica
~$\overline{\kres(s)}$ de~$\kres(s)$, un punto geométrico~$a_s$ de
\ifthenelse{\boolean{orig}}{$S_s$}{$X_s$}
con valores en esta extensión y una clase de caminos desde~$a_s$ hasta~$a$, de donde
se obtiene un homomorfismo~$\pi_1(\overline X_s,a_s)\to\pi_1(X,a)$, donde
~$\overline X_s=X_s\otimes_{\kres(s)}\overline{\kres(s)}$. Entonces el
homomorfismo~$\pi_1(X,a)\to\pi_1(S,b)$ es sobreyectivo, y su núcleo es el
subgrupo normal cerrado de~$\pi_1(X,a)$ engendrado por las imágenes de los
~$\pi_1(\overline X_s,a_s)$.
\end{corollaire}

\begin{remarque}
\label{IX.6.12}
Bajo las condiciones de~\Ref{IX.6.7}, suponiendo~$S$ noetheriano, se ve
fácilmente que el conjunto de los puntos~$s\in S$ tales que
\ifthenelse{\boolean{orig}}{$S'_{s'}$}{$X'_s$}
sea geométricamente trivial sobre~$X_s$ es constructible; si~$X'$ es propio
sobre~$X$, es incluso abierto, como se ve en~\Ref{IX.6.6}. Esto permite, por
tanto, si~$S$
\ifthenelse{\boolean{orig}}{y}{es}
un preesquema de Jacobson (por ejemplo, de tipo finito sobre un cuerpo), o
\marginpar{260}
si~$X'$ es propio sobre~$X$, limitarse, para verificar las condiciones
de~\Ref{IX.6.7}, a los puntos~$s$ de~$S$ que son \emph{cerrados}. Del mismo
modo, en~\Ref{IX.6.11} basta entonces tomar los
~$\pi_1(\overline X_s,a_s)$ para los puntos de~$S$ que son cerrados.
\end{remarque}

\begin{sgabibliography}{D}{IX.7}

\bibitem[D]{IX.D} J\ptbl \textsc{Giraud}, Méthode de la
\ifthenelse{\boolean{orig}}
{d'escente,}
{descente,}
Mémoire \no 2 de la Soc.\ Math.\ Fr.,~1964.

\bibitem[1]{IX.1} A\ptbl \textsc{Grothendieck}, Géométrie Algébrique et
Géométrie Formelle, Séminaire Bourbaki t\ptbl 11, 1959, \No 182.
\end{sgabibliography}
